\section{Introduction}
\label{sec:intro}

Star formation in galaxies is a multi-scale, multi-physics process whose
regulation has emerged as one of the central problems of galaxy evolution.
Across a wide range of galactic environments, observations reveal tight
correlations between local gas properties, dynamical equilibrium pressure,
and the star formation rate (SFR) surface density
\citep{2010ApJ...721..975O,2011ApJ...731...41O,2020ApJ...892..148S,2022ApJ...936..137O}.
These correlations are now understood as a signature of self-regulation:
energy and momentum injected by massive-star feedback---supernovae, stellar
winds, and radiation pressure---maintain the interstellar medium (ISM) in
an approximately steady state in which the vertical weight of the gas is
supported against gravity by turbulent, thermal, and magnetic pressures.
The resulting balance ties the local SFR to the local dynamical state of
the disk, and reproduces the observed scaling between $\Ssfr$ and the
dynamical equilibrium pressure $\PDE$ across more than four orders of
magnitude in environment
\citep{2020ApJ...892..148S,2023ApJ...945L..19S,2025ApJ...989...66V}.

The pressure-regulated, feedback-modulated (\PRFM) framework
\citep{2022ApJ...936..137O} provides a theoretical foundation for this
self-regulation. In its compact form, vertical equilibrium fixes the
midplane total pressure to the weight of the overlying gas, $\Ptot = \W$,
where $\W$ is the integrated weight of the gas column in the combined
gravitational potential of the gas, stars, and dark matter. The feedback
yield $\Ytot \equiv \Ptot/\Ssfr$ then encapsulates how efficiently young
stellar populations convert star formation into the pressure required to
support the disk. Equating these two relations yields the prediction
\begin{equation}
  \Ssfr = \frac{\W}{\Ytot},
  \label{eq:prfm}
\end{equation}
in which $\W$ is dictated by the local galactic environment and $\Ytot$ is
set by feedback physics. The \TIGRESS\ simulation framework
\citep{2024ApJ...972...67K} resolves the multiphase ISM in local
shearing-box magnetohydrodynamic (MHD) simulations with explicit
non-equilibrium cooling/heating, UV radiation transfer, cosmic-ray
transport, and clustered supernova feedback, and delivers calibrated values
of $\Ytot$ and its thermal, turbulent, and magnetic components across
metallicities. \PRFM\ has been adopted as the star formation prescription
in cosmological-volume simulations \citep{2024ApJ...975..151H} and
isolated-galaxy simulations \citep{2024ApJ...975..113J}, where it provides
a physically motivated alternative to empirical Kennicutt--Schmidt-type
prescriptions.

The observational counterpart to \TIGRESS\ is provided by the Physics at
High Angular resolution in Nearby GalaxieS (\PHANGS) survey
\citep{2021ApJS..255...19L,2023AJ....166..145L}, which delivers
spatially resolved, multi-wavelength measurements of $\sim80$ nearby
star-forming disk galaxies. \PHANGS\ aperture catalogs combine high-resolution
CO mapping with matched \HI, stellar mass, and SFR tracers to produce
joint constraints on gas, stars, and star formation at sub-kpc scales
\citep{2022AJ....164...43S,2023ApJ...945L..19S}.
Together with stellar dynamical modeling of the same galaxies
\citep{2020ApJ...892..148S,2025ApJ...989...66V}, the survey defines an
empirical envelope of local galactic environments---spanning more than two
orders of magnitude in gas and stellar surface density, and the full range
of orbital frequencies found in disk-dominated systems---against which any
ISM theory must be tested.

To date, however, comparisons between \TIGRESS-type simulations and
\PHANGS-type observations have relied almost exclusively on a handful of
representative simulations selected by hand at the solar neighborhood and a
few extreme environments. No systematic survey has yet been carried out
that spans the full \PHANGS\ parameter volume with controlled, calibrated
simulations of the resolved multiphase ISM. An analogous gap in
cosmology---between hand-picked hydrodynamic simulations and the
high-dimensional space of cosmological and astrophysical parameters
required for inference---was closed by the Cosmology and Astrophysics with
MachinE Learning Simulations (\CAMELS) project
\citep{2021ApJ...915...71V}. \CAMELS\ demonstrated that a structured
simulation ensemble covering parameter space in a controlled way enables
emulator construction, simulation-based inference, and systematic
comparison with observations, even at the modest individual-simulation
fidelity that suite-scale projects can afford. The same logic motivates a
\CAMELS-style ensemble for the star-forming ISM.

In this paper, the first in a series, we present the design of a
\PHANGS-informed simulation suite for the star-forming ISM based on the
\TIGRESS\ framework. We use the \PHANGS\ megatable
\citep{2022AJ....164...43S} to construct an observational prior over the
five environmental parameters that enter \TIGRESS\ as initial conditions:
the gas surface density $\Sgas$, stellar surface density $\Sstar$,
stellar scale height $\Hstar$, orbital frequency $\Omrot$, and shear
parameter $\qshear$. We construct the simulation design by fitting a
kernel density estimator (KDE) to the joint \PHANGS\ distribution of these
five fields and mapping a scrambled Sobol quasi-random sequence through
its marginal quantile functions. The Sobol construction allows the suite
to be extended progressively---doubling the sample size without displacing
or re-running earlier simulations---and yields low-discrepancy coverage of
the \PHANGS\ environmental volume at every intermediate sample size. The
molecular fraction $\fmol$ and SFR surface density $\Ssfr$ are
deliberately excluded from the input design and reserved as validation
quantities to be compared against emergent simulation outputs. In this
first suite, all physics parameters (metallicity, dust-to-gas ratio,
feedback prescription) are held at fiducial solar-neighborhood values;
extension to physics-parameter variations and the use of this suite for
simulation-based inference (\SBI) are the subjects of companion papers
\CGK{(Paper~II, in prep.)}.

The remainder of this paper is organized as follows. \autoref{sec:data}
describes the \PHANGS\ megatable data, the aperture-level join, and the
derivation of the five design fields and validation fields used throughout.
\CGK{Section~3} presents the sampling design, including the KDE fit, the
Sobol construction, and the fairness diagnostics used to assess marginal
and joint coverage. \autoref{sec:tigress_mapping} describes the mapping
from \PHANGS\ design fields to \TIGRESS\ initial conditions and the
fiducial physics choices for the first suite. \autoref{sec:results}
presents the resulting simulation design, an initial set of \TIGRESS\ runs,
and validation against \PHANGS-observed $\fmol$, $\Ssfr$, $\PDE$, and
$\seff$. \CGK{Section~6} discusses extensions to physics-parameter
variations, the role of this suite in enabling \SBI\ on the resolved ISM,
and connections to large-volume cosmological simulations adopting \PRFM.
